\section*{Probability, Bayes, densities}
\subsection*{Conditional probability and Bayes}
\textbf{Purpose:} turn likelihoods and priors into class probabilities after seeing $x$.
\[
\fid[eq:bayes]\begin{aligned}
P(A|B)&=\frac{P(A,B)}{P(B)},\\
P(A)&=\sum_cP(A|c)P(c),\\
P(y=c|x)&=\frac{P(x|y=c)P(y=c)}
{\sum_{c'}P(x|y=c')P(y=c')}.
\end{aligned}
\]
\textbf{Symbols:} prior $P(y=c)$ is class frequency before $x$; likelihood $P(x|y=c)$ says how likely $x$ is inside class $c$; posterior $P(y=c|x)$ is the final class probability.
\textbf{Do this:} use \fref{eq:bayes}; for each class compute score = likelihood $\times$ prior; add all scores; divide each score by the sum; choose the largest posterior.

\subsection*{Counts and Naive Bayes}
\[
\fid[eq:naive-bayes]\begin{aligned}
\hat P(y=c|x=a)&=
\frac{\#\{i:y_i=c,x_i=a\}}{\#\{i:x_i=a\}},\\
\operatorname{score}(c)&=P(y=c)\prod_{m=1}^MP(x_m|y=c),\\
P(y=c|x)&=\frac{\operatorname{score}(c)}{\sum_{c'}\operatorname{score}(c')},\\
\log\operatorname{score}(c)&=\log P(c)+\sum_m\log P(x_m|c).
\end{aligned}
\]
\textbf{How to use tables:} for each class row/column, pick the probability for the observed value of each attribute; multiply them with the prior as in \fref{eq:naive-bayes}; normalize scores so posteriors sum to $1$. Use log-scores when products are tiny; the largest log-score gives the same class.
\textbf{From counts:} if a table gives raw counts, convert each needed entry to a conditional probability by dividing by the total count in the same condition/class. For $P(x_m=a|y=c)$, denominator is the number of rows in class $c$, not all rows.

\subsection*{Gaussian density classifier}
\[
\fid[eq:gauss]\begin{aligned}
\mathcal N(x|\mu,\Sigma)&=
\frac{\exp[-\tfrac12(x-\mu)^\top\Sigma^{-1}(x-\mu)]}
{(2\pi)^{M/2}|\Sigma|^{1/2}},\\
\mathcal N(x|\mu,\sigma^2)&=\frac1{\sqrt{2\pi}\sigma}
\exp\!\left[-\frac{(x-\mu)^2}{2\sigma^2}\right),\\
\hat y&=\argmax_c\{\mathcal N(x|\mu_c,\Sigma_c)P(c)\}.
\end{aligned}
\]
\textbf{Symbols:} $\mu$ is centre, $\Sigma$ covariance, $\sigma$ standard deviation, $|\Sigma|$ determinant. Density height is not a probability by itself.
\textbf{Do this:} use \fref{eq:gauss}; for each class insert $x,\mu_c,\Sigma_c$ or $\sigma_c$; multiply density by prior; choose the largest class score. Curves are highest near $\mu$; larger variance makes a wider, lower curve.
\textbf{Boundary reading:} with equal priors, choose the class with larger density at $x$; with unequal priors, a class with lower density can still win if its prior is large enough.
